\documentclass{article}

\begin{document}

\section{Bubble sort}
\subsection{Core concepts}
  This sorting algorithm compares two adjacent elements and swap their positions if they are in wrong order. As a result, the largest element will move to the end of the array. This movement is similar to bubbles rising to the surface, that is why this sorting method is called bubble sort.

\subsection{Step-by-step explanation}
  This algorithm will sort the array using several pass. After the first pass, the maximum element goes to the end (its correct position). Similarly, the second largest goes to the second last position and so on.

Example: the given array has 6 elements: 20, 80, 40, 25, 60, 30



The algorithm will iterate 5 times
\begin{itemize}
    \item Pass 1: 20 40 25 60 30 80 
    \item Pass 2: 20 25 40 30 60 80 
    \item Pass 3: 20 25 30 40 60 80 
    \item Pass 4: 20 25 30 40 60 80  
    \item Pass 5: 20 25 30 40 60 80  
\end{itemize}
 
 

\subsection{Complexity Analysis}

Comparisons:  
The inner loop executes the size of unsorted part minus 1, and in each iteration, there is one key comparison.

Assignment:  
For the worst case, the number of swaps equals to the number of comparisons, but each swap requires 3 assignments.

In the worst case, this algorithm requires

\[
(1+3)((n-1)+(n-2)+\dots+2+1)=\frac{4n(n-1)}{2}=2n(n-1)
\]

Big-O complexity of bubble sort is

\[
O(n^2)
\]

\subsection{Variants and optimizations}

This technique has another version which adds a boolean flag to check if the array is sorted. This will reduce the complexity to

\[
O(n)
\]

in the best case when the array is already sorted.

\end{document}